\section{Zadanie 16}
\autor{Mateusz Wojaczek}
Niech \(X\) i \(Y\) mają łączną gęstość \(f(x, y) = cx(y - x)e^{-y}\) dla \(0 \le x \le y < \infty\).

\begin{enumerate}
	\item Wyznacz wartość \(c\).
	\item Udowodnij, że
	      \begin{align*}
		      f_{X|Y}(x | y) & = 6x(y - x)y^{-3} \quad \text{dla } 0 \le x \le y,           \\
		      f_{Y|X}(y | x) & = (y - x)e^{x - y} \quad \text{dla } 0 \le x \le y < \infty.
	      \end{align*}
	\item Wynioskuj, że \(\text{E}(X | Y) = \frac{1}{2}Y\) oraz \(\text{E}(Y | X) = X + 2\).
\end{enumerate}

\subsection*{Rozwiązanie}
\subsection*{(i) Wyznaczenie stałej \(c\)}

Aby \(f_{X,Y}(x,y)\) była gęstością, całka podwójna musi wynosić 1:
\[
	\int_0^\infty \int_0^y c x (y-x) e^{-y} \, dx \, dy = 1.
\]

Najpierw całkujemy względem \(x\):
\[
	\int_0^y x(y-x) dx = \int_0^y (xy - x^2) dx = \frac{1}{2}y^3 - \frac{1}{3}y^3 = \frac{1}{6}y^3.
\]

Zatem:
\[
	\int_0^\infty c \cdot \frac{1}{6} y^3 e^{-y} dy = 1 \implies \frac{c}{6} \int_0^\infty y^3 e^{-y} dy = 1.
\]

\(\int_0^\infty y^3 e^{-y} dy = 6\), więc:
\[
	\frac{c}{6} \cdot 6 = c = 1.
\]

\textbf{Odpowiedź:} \(c = 1\).

\subsection*{(ii) Wyznaczenie gęstości warunkowych}

\paragraph{Gęstość warunkowa \(X\) przy \(Y=y\)}:

\[
	f_{X|Y}(x|y) = \frac{f_{X,Y}(x,y)}{f_Y(y)}.
\]

Najpierw obliczamy \(f_Y(y)\):
\[
	f_Y(y) = \int_0^y f_{X,Y}(x,y) dx = \int_0^y 6 x (y-x) e^{-y} dx = 6 e^{-y} \cdot \frac{1}{6} y^3 = y^3 e^{-y}.
\]

Zatem:
\[
	f_{X|Y}(x|y) = \frac{6 x (y-x) e^{-y}}{y^3 e^{-y}} = \frac{6 x (y-x)}{y^3}, \quad 0 \le x \le y.
\]

\paragraph{Gęstość warunkowa \(Y\) przy \(X=x\)}:

\[
	f_{Y|X}(y|x) = \frac{f_{X,Y}(x,y)}{f_X(x)}.
\]

Najpierw obliczamy \(f_X(x)\):
\[
	f_X(x) = \int_x^\infty 6 x (y-x) e^{-y} dy = 6x \int_x^\infty (y-x) e^{-y} dy.
\]

Podstawiamy \(u = y-x \implies dy = du, y-x = u, y = u+x\):
\[
	\int_x^\infty (y-x) e^{-y} dy = \int_0^\infty u e^{-(u+x)} du = e^{-x} \int_0^\infty u e^{-u} du = e^{-x} \cdot 1! = e^{-x}.
\]

Zatem \(f_X(x) = 6 x e^{-x}\). Wobec tego:
\[
	f_{Y|X}(y|x) = \frac{6 x (y-x) e^{-y}}{6 x e^{-x}} = (y-x) e^{-(y-x)}, \quad y \ge x.
\]

\subsection*{(iii) Wartości oczekiwane warunkowe}

\paragraph{Dla \(E(X|Y)\)}:

\[
	E(X|Y=y) = \int_0^y x f_{X|Y}(x|y) dx = \int_0^y x \cdot \frac{6x(y-x)}{y^3} dx = \frac{6}{y^3} \int_0^y (x^2 y - x^3) dx.
\]

\[
	\int_0^y (x^2 y - x^3) dx = y \cdot \frac{y^3}{3} - \frac{y^4}{4} = \frac{y^4}{3} - \frac{y^4}{4} = \frac{y^4}{12}.
\]

\[
	E(X|Y=y) = \frac{6}{y^3} \cdot \frac{y^4}{12} = \frac{y}{2}.
\]

\paragraph{Dla \(E(Y|X)\)}:

\[
	E(Y|X=x) = \int_x^\infty y f_{Y|X}(y|x) dy = \int_x^\infty y (y-x) e^{-(y-x)} dy.
\]

Podstawiamy \(u = y-x \implies y = u+x, dy = du\):
\[
	E(Y|X=x) = \int_0^\infty (u+x) u e^{-u} du = \int_0^\infty u^2 e^{-u} du + x \int_0^\infty u e^{-u} du = 2 + x \cdot 1 = x + 2.
\]


\[
	E(X|Y) = \frac{1}{2}Y, \quad E(Y|X) = X+2.
\]

\subsection*{Odpowiedź}
\[
	\boxed{
		\begin{aligned}
			 & \textbf{Wartość } c\textbf{:}           &  & \mathbf{c = 1}                                                                       \\
			 & \textbf{Gęstości warunkowe:}            &  & f_{X|Y}(x | y) = \mathbf{6x(y - x)y^{-3}}, \quad \text{dla } 0 \le x \le y           \\
			 &                                         &  & f_{Y|X}(y | x) = \mathbf{(y - x)e^{x - y}}, \quad \text{dla } 0 \le x \le y < \infty \\
			 & \textbf{Wartości oczekiwane warunkowe:} &  & \mathbf{E(X | Y) = \frac{1}{2}Y}                                                     \\
			 &                                         &  & \mathbf{E(Y | X) = X + 2}
		\end{aligned}
	}
\]
